\documentclass[a4paper,11pt]{article}

\usepackage[margin=2.5cm]{geometry}
\usepackage{amsmath, amssymb, amsthm}
\usepackage{physics}
\usepackage{bm}
\usepackage{siunitx}
\usepackage{booktabs}
\usepackage{longtable}
\usepackage{hyperref}

\setlength{\tabcolsep}{10pt}
\renewcommand{\arraystretch}{1.6}

\title{Technical Note: Design Theory for the Butterworth-Optimized\\
Photodiode Transimpedance Amplifier}
\author{N.~Moteki}
\date{Last updated: 2026-03-19 (corresponds to Photodiode-TIA v0.1.0)}

\begin{document}
\maketitle

\begin{abstract}
This document provides a self-contained description of the mathematical
theory implemented in the photodiode transimpedance amplifier (TIA)
design tool.  It covers the single-pole opamp model, the derivation of
the closed-loop second-order transfer function, the exact
Butterworth-optimum feedback capacitance without the common
$\omega_{\mathrm{c}} R_{\mathrm{f}} C_{\mathrm{f}} \gg 1$
approximation, the $-3\,\mathrm{dB}$ bandwidth, the noise equivalent
bandwidth for a second-order Butterworth response, the piecewise
noise-gain integration for opamp voltage noise, and the shot/thermal
noise analysis.  All algebraic steps are shown explicitly.
\end{abstract}

\tableofcontents
\newpage

%======================================================================
\section{Circuit model and assumptions}
\label{sec:circuit-model}
%======================================================================

\subsection{Circuit topology}

The TIA is an inverting opamp configuration with a photodiode at the
input.  The non-inverting input is connected to ground.

\begin{itemize}
  \item $I_{\mathrm{pd}}$: photocurrent from the photodiode (signal
    source, modelled as ideal current source)
  \item $R_{\mathrm{f}}$, $C_{\mathrm{f}}$: feedback resistance and
    capacitance (in parallel), connected from the output to the
    inverting input
  \item $C_{\mathrm{i}}$: total input capacitance at the inverting
    node,
    $C_{\mathrm{i}} = C_{\mathrm{d}} + C_{\mathrm{icm}} + C_{\mathrm{id}}$
    \begin{itemize}
      \item $C_{\mathrm{d}}$: photodiode junction capacitance
      \item $C_{\mathrm{icm}}$: opamp common-mode input capacitance
      \item $C_{\mathrm{id}}$: opamp differential input capacitance
    \end{itemize}
  \item $C_{\mathrm{s}}$: stray (parasitic) capacitance at the
    inverting node to ground
\end{itemize}

\subsection{Opamp model}

The opamp is modelled with a single-pole open-loop gain:
\begin{equation}
  A(s) = \frac{A_0 \,\omega_{\mathrm{p}}}{s + \omega_{\mathrm{p}}}
  \label{eq:opamp-full}
\end{equation}
where $A_0$ is the DC open-loop gain and $\omega_{\mathrm{p}}$ is the
dominant pole frequency.  For frequencies well above the dominant pole
($s \gg \omega_{\mathrm{p}}$, i.e.\ $f \gg f_{\mathrm{p}}$, typically
satisfied for $f > \SI{10}{\hertz}$), this simplifies to:
\begin{equation}
  A(s) \approx \frac{A_0 \,\omega_{\mathrm{p}}}{s}
  = \frac{\omega_{\mathrm{c}}}{s}
  \label{eq:opamp-integrator}
\end{equation}
where $\omega_{\mathrm{c}} = A_0\,\omega_{\mathrm{p}} = 2\pi
f_{\mathrm{c}}$ is the gain-bandwidth product (GBW) in
$\si{\radian\per\second}$.  This ``integrator model'' is the standard
approximation used throughout TIA design literature and is accurate
across the entire frequency range of interest
($\si{\kilo\hertz}$--$\si{\mega\hertz}$).

%======================================================================
\section{Closed-loop transimpedance transfer function}
\label{sec:transfer-function}
%======================================================================

\subsection{Goal}

Derive the closed-loop transimpedance:
\begin{equation*}
  Z_{\mathrm{TI}}(s)
  \equiv \frac{V_{\mathrm{out}}(s)}{I_{\mathrm{pd}}(s)}.
\end{equation*}

\subsection{Step 1: Define the feedback impedance}

$R_{\mathrm{f}}$ and $C_{\mathrm{f}}$ are in parallel:
\begin{equation}
  Z_{\mathrm{f}}
  = R_{\mathrm{f}} \;\|\; \frac{1}{sC_{\mathrm{f}}}
  = \frac{R_{\mathrm{f}}}{1 + sR_{\mathrm{f}}C_{\mathrm{f}}}
  \label{eq:Zf}
\end{equation}

\subsection{Step 2: Opamp constraint}

With the non-inverting input grounded:
\begin{equation*}
  V_{\mathrm{out}} = -A(s) \cdot V^{-}
  = -\frac{\omega_{\mathrm{c}}}{s}\,V^{-}
\end{equation*}
Solving for $V^-$:
\begin{equation}
  V^{-} = -\frac{s}{\omega_{\mathrm{c}}}\,V_{\mathrm{out}}
  \label{eq:Vminus}
\end{equation}

\subsection{Step 3: Kirchhoff's current law at the inverting node}

All currents leaving the inverting node must sum to zero.  The current
into the node from the photodiode is $I_{\mathrm{pd}}$.  The currents
leaving the node are:
\begin{itemize}
  \item Through the feedback network to the output:
    $(V^{-} - V_{\mathrm{out}}) / Z_{\mathrm{f}}$
  \item Through the total shunt capacitance
    $(C_{\mathrm{i}} + C_{\mathrm{s}})$ to ground:
    $V^{-} \cdot s(C_{\mathrm{i}} + C_{\mathrm{s}})$
\end{itemize}
KCL:
\begin{equation*}
  I_{\mathrm{pd}}
  + \frac{V^{-} - V_{\mathrm{out}}}{Z_{\mathrm{f}}}
  + V^{-} \cdot s(C_{\mathrm{i}} + C_{\mathrm{s}}) = 0
\end{equation*}
Note the sign convention: current flowing from the node through
$Z_{\mathrm{f}}$ toward $V_{\mathrm{out}}$ is
$(V^{-} - V_{\mathrm{out}})/Z_{\mathrm{f}}$.  Rearranging with
current into the node on the left:
\begin{equation}
  I_{\mathrm{pd}}
  = \frac{V_{\mathrm{out}} - V^{-}}{Z_{\mathrm{f}}}
  - V^{-} \cdot s(C_{\mathrm{i}} + C_{\mathrm{s}})
  \label{eq:kcl}
\end{equation}

\subsection{Step 4: Substitute the opamp constraint}

Substitute \eqref{eq:Zf} and \eqref{eq:Vminus} into \eqref{eq:kcl}.

First, compute $V_{\mathrm{out}} - V^-$:
\begin{equation}
  V_{\mathrm{out}} - V^{-}
  = V_{\mathrm{out}}\!\left(1 + \frac{s}{\omega_{\mathrm{c}}}\right)
  = V_{\mathrm{out}} \cdot
    \frac{\omega_{\mathrm{c}} + s}{\omega_{\mathrm{c}}}
  \label{eq:Vdiff}
\end{equation}

Substitute into the feedback current term:
\begin{equation}
  \frac{V_{\mathrm{out}} - V^{-}}{Z_{\mathrm{f}}}
  = V_{\mathrm{out}} \cdot
    \frac{\omega_{\mathrm{c}} + s}{\omega_{\mathrm{c}}} \cdot
    \frac{1 + sR_{\mathrm{f}}C_{\mathrm{f}}}{R_{\mathrm{f}}}
  \label{eq:Ifb}
\end{equation}

The shunt capacitance current term:
\begin{equation}
  -V^{-} \cdot s(C_{\mathrm{i}} + C_{\mathrm{s}})
  = \frac{s^2(C_{\mathrm{i}} + C_{\mathrm{s}})}{\omega_{\mathrm{c}}}
    \,V_{\mathrm{out}}
  \label{eq:Ishunt}
\end{equation}

\subsection{Step 5: Combine}

Substituting \eqref{eq:Ifb} and \eqref{eq:Ishunt} into
\eqref{eq:kcl}:
\begin{equation*}
  I_{\mathrm{pd}}
  = V_{\mathrm{out}} \left[
    \frac{(\omega_{\mathrm{c}} + s)
          (1 + sR_{\mathrm{f}}C_{\mathrm{f}})}
         {\omega_{\mathrm{c}} R_{\mathrm{f}}}
    + \frac{s^2(C_{\mathrm{i}} + C_{\mathrm{s}})}{\omega_{\mathrm{c}}}
  \right]
\end{equation*}
Factor out $1/\omega_{\mathrm{c}}$:
\begin{equation}
  I_{\mathrm{pd}}
  = \frac{V_{\mathrm{out}}}{\omega_{\mathrm{c}}} \left[
    \frac{(\omega_{\mathrm{c}} + s)
          (1 + sR_{\mathrm{f}}C_{\mathrm{f}})}
         {R_{\mathrm{f}}}
    + s^2(C_{\mathrm{i}} + C_{\mathrm{s}})
  \right]
  \label{eq:Ipd-combined}
\end{equation}

\subsection{Step 6: Expand the product}

\begin{align*}
  (\omega_{\mathrm{c}} + s)(1 + sR_{\mathrm{f}}C_{\mathrm{f}})
  &= \omega_{\mathrm{c}}
   + \omega_{\mathrm{c}} s R_{\mathrm{f}} C_{\mathrm{f}}
   + s + s^2 R_{\mathrm{f}} C_{\mathrm{f}} \\
  &= s^2 R_{\mathrm{f}} C_{\mathrm{f}}
   + s(1 + \omega_{\mathrm{c}} R_{\mathrm{f}} C_{\mathrm{f}})
   + \omega_{\mathrm{c}}
\end{align*}

\subsection{Step 7: Divide by $R_{\mathrm{f}}$ and add the shunt term}

\begin{align*}
  &\frac{(\omega_{\mathrm{c}} + s)
        (1 + sR_{\mathrm{f}}C_{\mathrm{f}})}
       {R_{\mathrm{f}}}
  + s^2(C_{\mathrm{i}} + C_{\mathrm{s}}) \\
  &= s^2 C_{\mathrm{f}}
   + \frac{s(1 + \omega_{\mathrm{c}} R_{\mathrm{f}} C_{\mathrm{f}})}
          {R_{\mathrm{f}}}
   + \frac{\omega_{\mathrm{c}}}{R_{\mathrm{f}}}
   + s^2(C_{\mathrm{i}} + C_{\mathrm{s}}) \\
  &= s^2(C_{\mathrm{i}} + C_{\mathrm{f}} + C_{\mathrm{s}})
   + s\,\frac{1 + \omega_{\mathrm{c}} R_{\mathrm{f}} C_{\mathrm{f}}}
             {R_{\mathrm{f}}}
   + \frac{\omega_{\mathrm{c}}}{R_{\mathrm{f}}}
\end{align*}

\subsection{Step 8: Form the transfer function}

From \eqref{eq:Ipd-combined} and the result of Step~7:
\begin{equation}
  Z_{\mathrm{TI}}(s)
  = \frac{V_{\mathrm{out}}}{I_{\mathrm{pd}}}
  = \frac{\omega_{\mathrm{c}}}
         {s^2(C_{\mathrm{i}} + C_{\mathrm{f}} + C_{\mathrm{s}})
          + s\,\dfrac{1 + \omega_{\mathrm{c}}
          R_{\mathrm{f}} C_{\mathrm{f}}}{R_{\mathrm{f}}}
          + \dfrac{\omega_{\mathrm{c}}}{R_{\mathrm{f}}}}
  \label{eq:Zti-raw}
\end{equation}
(The overall sign is negative for the inverting configuration; we write
the magnitude here and note the $\SI{180}{\degree}$ phase inversion
separately.)

\subsection{Step 9: Identify the standard second-order form}

A standard second-order low-pass transfer function has the form:
\begin{equation*}
  H(s) = \frac{K\,\omega_{\mathrm{n}}^2}
              {s^2 + 2\zeta\omega_{\mathrm{n}}\,s
               + \omega_{\mathrm{n}}^2}
\end{equation*}

Dividing numerator and denominator of \eqref{eq:Zti-raw} by
$(C_{\mathrm{i}} + C_{\mathrm{f}} + C_{\mathrm{s}})$:
\begin{equation}
  Z_{\mathrm{TI}}(s)
  = \frac{\dfrac{\omega_{\mathrm{c}}}
              {C_{\mathrm{i}} + C_{\mathrm{f}} + C_{\mathrm{s}}}}
         {s^2
          + \dfrac{1 + \omega_{\mathrm{c}}R_{\mathrm{f}}C_{\mathrm{f}}}
                  {R_{\mathrm{f}}
                   (C_{\mathrm{i}} + C_{\mathrm{f}} + C_{\mathrm{s}})}\,s
          + \dfrac{\omega_{\mathrm{c}}}
                  {R_{\mathrm{f}}
                   (C_{\mathrm{i}} + C_{\mathrm{f}} + C_{\mathrm{s}})}}
  \label{eq:Zti-standard}
\end{equation}

By comparison with the standard form, we identify:
\begin{equation}
  \omega_{\mathrm{n}}^2
  = \frac{\omega_{\mathrm{c}}}
         {R_{\mathrm{f}}(C_{\mathrm{i}} + C_{\mathrm{f}}
          + C_{\mathrm{s}})}
  \label{eq:omega-n}
\end{equation}
\begin{equation}
  2\zeta\omega_{\mathrm{n}}
  = \frac{1 + \omega_{\mathrm{c}} R_{\mathrm{f}} C_{\mathrm{f}}}
         {R_{\mathrm{f}}(C_{\mathrm{i}} + C_{\mathrm{f}}
          + C_{\mathrm{s}})}
  \label{eq:damping}
\end{equation}

The DC gain ($s = 0$) of \eqref{eq:Zti-standard}:
\begin{equation*}
  Z_{\mathrm{TI}}(0)
  = \frac{\omega_{\mathrm{c}}
          / (C_{\mathrm{i}} + C_{\mathrm{f}} + C_{\mathrm{s}})}
         {\omega_{\mathrm{c}}
          / [R_{\mathrm{f}}(C_{\mathrm{i}} + C_{\mathrm{f}}
          + C_{\mathrm{s}})]}
  = R_{\mathrm{f}}.
\end{equation*}

The numerator constant can therefore be written as
$R_{\mathrm{f}}\,\omega_{\mathrm{n}}^2$, giving the final form
(restoring the sign for the inverting configuration):
\begin{equation}
  Z_{\mathrm{TI}}(s)
  = \frac{-R_{\mathrm{f}}\,\omega_{\mathrm{n}}^2}
         {s^2 + 2\zeta\omega_{\mathrm{n}}\,s
          + \omega_{\mathrm{n}}^2}
  \label{eq:Zti-final}
\end{equation}

%======================================================================
\section{Butterworth-optimum feedback capacitance}
\label{sec:butterworth}
%======================================================================

\subsection{Objective}

Find the value of $C_{\mathrm{f}}$ that makes the closed-loop response
maximally flat (Butterworth).

\subsection{Step 1: Butterworth condition}

A second-order Butterworth (maximally flat magnitude) filter has
quality factor:
\begin{equation*}
  Q = \frac{1}{\sqrt{2}}
\end{equation*}
The relationship between $Q$ and the damping ratio $\zeta$ is
$Q = 1/(2\zeta)$, so:
\begin{equation}
  \zeta = \frac{1}{\sqrt{2}}
  \label{eq:zeta-bw}
\end{equation}
This is the unique damping ratio that produces zero gain peaking in the
magnitude response (i.e.\ $|Z_{\mathrm{TI}}(j\omega)|$ is
monotonically decreasing).

\subsection{Step 2: Translate the condition to circuit parameters}

The Butterworth condition \eqref{eq:zeta-bw} implies:
\begin{equation*}
  (2\zeta\omega_{\mathrm{n}})^2
  = 4\zeta^2\omega_{\mathrm{n}}^2
  = 4 \cdot \frac{1}{2} \cdot \omega_{\mathrm{n}}^2
  = 2\,\omega_{\mathrm{n}}^2.
\end{equation*}

Substituting \eqref{eq:omega-n} and \eqref{eq:damping}:
\begin{equation*}
  \left(
    \frac{1 + \omega_{\mathrm{c}} R_{\mathrm{f}} C_{\mathrm{f}}}
         {R_{\mathrm{f}}(C_{\mathrm{i}} + C_{\mathrm{f}}
          + C_{\mathrm{s}})}
  \right)^{\!2}
  = 2 \cdot \frac{\omega_{\mathrm{c}}}
                  {R_{\mathrm{f}}(C_{\mathrm{i}} + C_{\mathrm{f}}
                   + C_{\mathrm{s}})}
\end{equation*}
Multiply both sides by
$[R_{\mathrm{f}}(C_{\mathrm{i}} + C_{\mathrm{f}} + C_{\mathrm{s}})]^2$:
\begin{equation}
  (1 + \omega_{\mathrm{c}} R_{\mathrm{f}} C_{\mathrm{f}})^2
  = 2\,\omega_{\mathrm{c}} R_{\mathrm{f}}
    (C_{\mathrm{i}} + C_{\mathrm{f}} + C_{\mathrm{s}})
  \label{eq:bw-condition}
\end{equation}

\subsection{Step 3: Solve for $C_{\mathrm{f}}$}

Define a characteristic capacitance:
\begin{equation}
  C_{\mathrm{c}}
  \equiv \frac{1}{\omega_{\mathrm{c}} R_{\mathrm{f}}}
  = \frac{1}{2\pi f_{\mathrm{c}} R_{\mathrm{f}}}
  \label{eq:Cc}
\end{equation}
This has a physical interpretation: $C_{\mathrm{c}}$ is the capacitance
whose impedance at $f_{\mathrm{c}}$ equals $R_{\mathrm{f}}$.

Rewrite \eqref{eq:bw-condition} using
$\omega_{\mathrm{c}} R_{\mathrm{f}} = 1/C_{\mathrm{c}}$:
\begin{equation*}
  \left(1 + \frac{C_{\mathrm{f}}}{C_{\mathrm{c}}}\right)^{\!2}
  = \frac{2(C_{\mathrm{i}} + C_{\mathrm{f}} + C_{\mathrm{s}})}
         {C_{\mathrm{c}}}
\end{equation*}
Expand the left side:
\begin{equation*}
  1 + \frac{2C_{\mathrm{f}}}{C_{\mathrm{c}}}
    + \frac{C_{\mathrm{f}}^2}{C_{\mathrm{c}}^2}
  = \frac{2(C_{\mathrm{i}} + C_{\mathrm{s}})}{C_{\mathrm{c}}}
    + \frac{2C_{\mathrm{f}}}{C_{\mathrm{c}}}
\end{equation*}
The $2C_{\mathrm{f}}/C_{\mathrm{c}}$ terms cancel:
\begin{equation*}
  1 + \frac{C_{\mathrm{f}}^2}{C_{\mathrm{c}}^2}
  = \frac{2(C_{\mathrm{i}} + C_{\mathrm{s}})}{C_{\mathrm{c}}}
\end{equation*}
Multiply through by $C_{\mathrm{c}}^2$:
\begin{equation*}
  C_{\mathrm{f}}^2
  = 2(C_{\mathrm{i}} + C_{\mathrm{s}})\,C_{\mathrm{c}}
    - C_{\mathrm{c}}^2
  = C_{\mathrm{c}}\bigl(2(C_{\mathrm{i}} + C_{\mathrm{s}})
    - C_{\mathrm{c}}\bigr)
\end{equation*}
Taking the positive root:
\begin{equation}
  C_{\mathrm{f}}
  = \sqrt{C_{\mathrm{c}}\bigl(2(C_{\mathrm{i}} + C_{\mathrm{s}})
    - C_{\mathrm{c}}\bigr)},
  \qquad
  C_{\mathrm{c}} = \frac{1}{2\pi f_{\mathrm{c}} R_{\mathrm{f}}}
  \label{eq:Cf-exact}
\end{equation}

\subsection{Validity condition}

\eqref{eq:Cf-exact} requires the argument of the square root to be
non-negative:
\begin{equation*}
  2(C_{\mathrm{i}} + C_{\mathrm{s}}) - C_{\mathrm{c}} > 0
  \quad\Longleftrightarrow\quad
  C_{\mathrm{i}} + C_{\mathrm{s}}
  > \frac{C_{\mathrm{c}}}{2}
  = \frac{1}{2\omega_{\mathrm{c}} R_{\mathrm{f}}}.
\end{equation*}
When this condition is violated, the opamp GBW is too low (or
$R_{\mathrm{f}}$ too small) to achieve a Butterworth response---the
circuit is fundamentally underdamped regardless of $C_{\mathrm{f}}$.

\subsection{Relation to the common approximation}

Many references simplify by assuming
$\omega_{\mathrm{c}} R_{\mathrm{f}} C_{\mathrm{f}} \gg 1$
(equivalently $C_{\mathrm{f}} \gg C_{\mathrm{c}}$).  Under this
approximation, \eqref{eq:bw-condition} becomes:
\begin{equation*}
  (\omega_{\mathrm{c}} R_{\mathrm{f}} C_{\mathrm{f}})^2
  \approx 2\,\omega_{\mathrm{c}} R_{\mathrm{f}}
    (C_{\mathrm{i}} + C_{\mathrm{s}})
  \quad\Longrightarrow\quad
  C_{\mathrm{f}}
  \approx \sqrt{\frac{2(C_{\mathrm{i}} + C_{\mathrm{s}})}
                     {\omega_{\mathrm{c}} R_{\mathrm{f}}}}
  = \sqrt{2(C_{\mathrm{i}} + C_{\mathrm{s}})\,C_{\mathrm{c}}}.
\end{equation*}
\eqref{eq:Cf-exact} is the \textbf{exact} solution that does not rely
on this approximation.  The difference is significant when
$C_{\mathrm{c}}$ is not negligible compared to
$C_{\mathrm{i}} + C_{\mathrm{s}}$ (i.e.\ when $R_{\mathrm{f}}$ is
small or $f_{\mathrm{c}}$ is low).

%======================================================================
\section{Bandwidth}
\label{sec:bandwidth}
%======================================================================

\subsection{The $-3\,\mathrm{dB}$ bandwidth of a second-order
  Butterworth}

For a general second-order transfer function with damping ratio
$\zeta$, the squared magnitude response is:
\begin{equation*}
  |H(j\omega)|^2
  = \frac{\omega_{\mathrm{n}}^4}
         {(\omega_{\mathrm{n}}^2 - \omega^2)^2
          + (2\zeta\omega_{\mathrm{n}}\omega)^2}
\end{equation*}
For the Butterworth case $\zeta = 1/\sqrt{2}$, substituting and
simplifying:
\begin{equation*}
  |H(j\omega)|^2
  = \frac{\omega_{\mathrm{n}}^4}
         {(\omega_{\mathrm{n}}^2 - \omega^2)^2
          + 2\omega_{\mathrm{n}}^2\omega^2}
\end{equation*}
Expand the denominator:
\begin{equation*}
  = \omega_{\mathrm{n}}^4
    - 2\omega_{\mathrm{n}}^2\omega^2 + \omega^4
    + 2\omega_{\mathrm{n}}^2\omega^2
  = \omega_{\mathrm{n}}^4 + \omega^4
\end{equation*}
Therefore:
\begin{equation}
  |H(j\omega)|^2
  = \frac{\omega_{\mathrm{n}}^4}
         {\omega_{\mathrm{n}}^4 + \omega^4}
  = \frac{1}{1 + (\omega/\omega_{\mathrm{n}})^4}
  \label{eq:bw-mag}
\end{equation}
This is the defining form of a second-order Butterworth response.

At $\omega = \omega_{\mathrm{n}}$:
\begin{equation*}
  |H(j\omega_{\mathrm{n}})|^2 = \frac{1}{2}
  \quad\Longrightarrow\quad
  |H(j\omega_{\mathrm{n}})| = \frac{1}{\sqrt{2}}
  \;(-3\,\mathrm{dB}).
\end{equation*}
Therefore the $-3\,\mathrm{dB}$ bandwidth is exactly the natural
frequency:
\begin{equation*}
  \omega_{-3\,\mathrm{dB}} = \omega_{\mathrm{n}}.
\end{equation*}
In terms of ordinary frequency:
\begin{equation}
  f_{-3\,\mathrm{dB}} = \frac{\omega_{\mathrm{n}}}{2\pi}
  \label{eq:f3dB}
\end{equation}

\subsection{Express in terms of circuit parameters}

From \eqref{eq:omega-n}:
\begin{equation*}
  \omega_{\mathrm{n}}
  = \sqrt{\frac{\omega_{\mathrm{c}}}
               {R_{\mathrm{f}}(C_{\mathrm{i}} + C_{\mathrm{f}}
                + C_{\mathrm{s}})}}
  = \sqrt{\frac{2\pi f_{\mathrm{c}}}
               {R_{\mathrm{f}}(C_{\mathrm{i}} + C_{\mathrm{f}}
                + C_{\mathrm{s}})}}
\end{equation*}
Therefore:
\begin{equation}
  f_{-3\,\mathrm{dB}}
  = \frac{\omega_{\mathrm{n}}}{2\pi}
  = \sqrt{\frac{f_{\mathrm{c}}}
               {2\pi R_{\mathrm{f}}(C_{\mathrm{i}} + C_{\mathrm{f}}
                + C_{\mathrm{s}})}}
  \label{eq:f3dB-circuit}
\end{equation}

%======================================================================
\section{Noise equivalent bandwidth}
\label{sec:neb}
%======================================================================

\subsection{Definition}

The noise equivalent bandwidth (NEB) of a filter $H(f)$ is the
bandwidth of an ideal rectangular (brick-wall) filter with the same
peak gain that passes the same total noise power:
\begin{equation*}
  \mathrm{NEB}
  = \frac{1}{|H_{\mathrm{max}}|^2}
    \int_0^{\infty} |H(f)|^2 \dd{f}
\end{equation*}
For a Butterworth filter, $|H_{\mathrm{max}}| = |H(0)| = 1$, so:
\begin{equation}
  \mathrm{NEB} = \int_0^{\infty} |H(f)|^2 \dd{f}
  \label{eq:neb-def}
\end{equation}

\subsection{General $n$-th order Butterworth NEB}

For an $n$-th order Butterworth filter, the squared magnitude response
is:
\begin{equation*}
  |H(f)|^2 = \frac{1}{1 + (f/f_{-3\,\mathrm{dB}})^{2n}}
\end{equation*}
Substituting into \eqref{eq:neb-def} with
$u = f/f_{-3\,\mathrm{dB}}$:
\begin{equation}
  \mathrm{NEB}
  = f_{-3\,\mathrm{dB}} \int_0^{\infty} \frac{\dd{u}}{1 + u^{2n}}
  \label{eq:neb-integral}
\end{equation}

The integral
$\displaystyle I_n = \int_0^{\infty} \frac{\dd{u}}{1 + u^{2n}}$ is a
standard result that can be evaluated using the substitution
$t = u^{2n}$ and recognizing the beta function.

Setting $t = u^{2n}$, so $u = t^{1/(2n)}$,
$\dd{u} = \frac{1}{2n}\,t^{1/(2n)-1}\dd{t}$:
\begin{equation*}
  I_n = \int_0^{\infty} \frac{1}{1+t} \cdot
        \frac{1}{2n}\,t^{1/(2n)-1}\dd{t}
      = \frac{1}{2n}\,
        B\!\left(\frac{1}{2n},\;1 - \frac{1}{2n}\right)
\end{equation*}
where $B(a,b) = \Gamma(a)\Gamma(b)/\Gamma(a+b)$ is the beta function.

Using the reflection formula
$\Gamma(z)\,\Gamma(1-z) = \pi/\sin(\pi z)$ with $z = 1/(2n)$:
\begin{equation}
  I_n = \frac{1}{2n} \cdot
        \frac{\pi}{\sin\!\left(\dfrac{\pi}{2n}\right)}
  \label{eq:In}
\end{equation}
Therefore:
\begin{equation}
  \mathrm{NEB}
  = f_{-3\,\mathrm{dB}} \cdot
    \frac{\pi}{2n\,\sin\!\left(\dfrac{\pi}{2n}\right)}
  \label{eq:neb-general}
\end{equation}

\subsection{Specialization to $n = 2$}

For the second-order Butterworth ($n = 2$):
\begin{equation*}
  \mathrm{NEB}
  = f_{-3\,\mathrm{dB}} \cdot
    \frac{\pi}{4 \cdot \dfrac{\sqrt{2}}{2}}
  = f_{-3\,\mathrm{dB}} \cdot \frac{\pi}{2\sqrt{2}}
  = f_{-3\,\mathrm{dB}} \cdot \frac{\pi\sqrt{2}}{4}
\end{equation*}
\begin{equation}
  \mathrm{NEB}
  = f_{-3\,\mathrm{dB}} \cdot \frac{\pi\sqrt{2}}{4}
  \approx 1.111\,f_{-3\,\mathrm{dB}}
  \label{eq:neb-2nd}
\end{equation}

\subsection{Comparison with the single-pole approximation}

For a single-pole (first-order) system ($n = 1$):
\begin{equation*}
  \mathrm{NEB}_{n=1}
  = f_{-3\,\mathrm{dB}} \cdot \frac{\pi}{2\sin(\pi/2)}
  = f_{-3\,\mathrm{dB}} \cdot \frac{\pi}{2}
  \approx 1.571\,f_{-3\,\mathrm{dB}}.
\end{equation*}
Many TIA references use this single-pole approximation even when the
circuit has a second-order response.  The second-order Butterworth NEB
is $1.111/1.571 = 0.707$ times the single-pole value---a 29\%
difference that directly affects computed noise levels.

%======================================================================
\section{Noise analysis}
\label{sec:noise}
%======================================================================

The TIA output noise is the root-sum-square (RSS) of several
independent noise sources.  Each source has a spectral density that is
shaped by the closed-loop gain before appearing at the output.  The
noise sources fall into two categories: (a)~voltage noise of the opamp,
which is amplified by the frequency-dependent noise gain, and
(b)~current/thermal noise sources whose spectra are flat and are shaped
only by the closed-loop transimpedance $|Z_{\mathrm{TI}}(f)|^2$.

\subsection{Noise gain}

The noise gain $G_{\mathrm{n}}(f)$ is the closed-loop gain from the
opamp's input voltage noise source $e_{\mathrm{ni}}$ to the output.
For a TIA with feedback impedance
$Z_{\mathrm{f}} = R_{\mathrm{f}} \| (1/sC_{\mathrm{f}})$ and total
input-to-ground capacitance $(C_{\mathrm{i}} + C_{\mathrm{s}})$, the
noise gain is:
\begin{equation*}
  G_{\mathrm{n}}(f)
  = 1 + \frac{Z_{\mathrm{in}}}{Z_{\mathrm{f}}}
  = 1 + \frac{C_{\mathrm{i}} + C_{\mathrm{s}}}
             {C_{\mathrm{f}} + C_{\mathrm{s}}}
    \cdot
    \frac{1 + s R_{\mathrm{f}} (C_{\mathrm{f}} + C_{\mathrm{s}})}
         {1 + s R_{\mathrm{f}} (C_{\mathrm{i}} + C_{\mathrm{f}}
          + C_{\mathrm{s}})}
\end{equation*}

This gain exhibits four distinct frequency regions that are exploited
for piecewise integration:

\begin{enumerate}
  \item \textbf{Low-frequency flat region}
    ($f < f_{\mathrm{zf}}$):
    $G_{\mathrm{n}} \approx 1$
  \item \textbf{Rising region}
    ($f_{\mathrm{zf}} < f < f_{\mathrm{pf}}$):
    $G_{\mathrm{n}} \approx f / f_{\mathrm{zf}}$
  \item \textbf{Plateau region}
    ($f_{\mathrm{pf}} < f < f_{\mathrm{i}}$):
    $G_{\mathrm{n}} \approx 1 + C_{\mathrm{i}} / (C_{\mathrm{f}}
    + C_{\mathrm{s}})$
  \item \textbf{Roll-off region}
    ($f > f_{\mathrm{i}}$):
    $G_{\mathrm{n}}$ decreases as the loop gain falls below unity
\end{enumerate}

where the characteristic frequencies are:
\begin{equation}
  f_{\mathrm{zf}}
  = \frac{1}{2\pi R_{\mathrm{f}}
    (C_{\mathrm{i}} + C_{\mathrm{f}} + C_{\mathrm{s}})}
  \label{eq:fzf}
\end{equation}
\begin{equation}
  f_{\mathrm{pf}}
  = \frac{1}{2\pi R_{\mathrm{f}}
    (C_{\mathrm{f}} + C_{\mathrm{s}})}
  \label{eq:fpf}
\end{equation}
\begin{equation}
  f_{\mathrm{i}}
  = f_{\mathrm{c}} \cdot
    \frac{C_{\mathrm{f}} + C_{\mathrm{s}}}
         {C_{\mathrm{i}} + C_{\mathrm{f}} + C_{\mathrm{s}}}
  \label{eq:fi}
\end{equation}
Here $f_{\mathrm{zf}}$ is the noise-gain zero frequency,
$f_{\mathrm{pf}}$ is the noise-gain pole frequency, and
$f_{\mathrm{i}}$ is the unity loop-gain (intersection) frequency.

\subsection{Opamp voltage noise (piecewise integration)}

The opamp input voltage noise spectral density has a $1/f$ component
below the corner frequency $f_{\mathrm{f}}$ and a white (flat) floor
$e_{\mathrm{nif}}$ above it:
\begin{equation}
  e_{\mathrm{ni}}(f)
  = e_{\mathrm{nif}} \sqrt{1 + f_{\mathrm{f}} / f}
  \label{eq:eni}
\end{equation}

The output-referred voltage noise RMS is obtained by integrating the
product of $e_{\mathrm{ni}}^2(f)$ and $G_{\mathrm{n}}^2(f)$ over
frequency, divided into five segments.  Defining a low-frequency cutoff
$f_1$ (set to $\SI{0.01}{\hertz}$ in the implementation):

\textbf{Region 1}---$1/f$ noise ($f_1$ to $f_{\mathrm{f}}$), noise
gain $\approx 1$:
\begin{equation}
  E_{\mathrm{noe},1}^2
  = e_{\mathrm{nif}}^2 \cdot f_{\mathrm{f}}
    \ln\!\left(\frac{f_{\mathrm{f}}}{f_1}\right)
  \label{eq:Enoe1}
\end{equation}

\textbf{Region 2}---white noise, flat gain ($f_{\mathrm{f}}$ to
$f_{\mathrm{zf}}$), noise gain $\approx 1$:
\begin{equation}
  E_{\mathrm{noe},2}^2
  = e_{\mathrm{nif}}^2 (f_{\mathrm{zf}} - f_{\mathrm{f}})
  \label{eq:Enoe2}
\end{equation}

\textbf{Region 3}---white noise, rising gain ($f_{\mathrm{zf}}$ to
$f_{\mathrm{pf}}$), noise gain $\approx f/f_{\mathrm{zf}}$:
\begin{equation}
  E_{\mathrm{noe},3}^2
  = \frac{e_{\mathrm{nif}}^2}{f_{\mathrm{zf}}^2} \cdot
    \frac{f_{\mathrm{pf}}^3 - f_{\mathrm{zf}}^3}{3}
  \label{eq:Enoe3}
\end{equation}

\textbf{Region 4}---white noise, plateau ($f_{\mathrm{pf}}$ to
$f_{\mathrm{i}}$), noise gain
$\approx 1 + C_{\mathrm{i}}/(C_{\mathrm{f}} + C_{\mathrm{s}})$:
\begin{equation}
  E_{\mathrm{noe},4}^2
  = e_{\mathrm{nif}}^2
    \left(1 + \frac{C_{\mathrm{i}}}
                   {C_{\mathrm{f}} + C_{\mathrm{s}}}\right)^{\!2}
    (f_{\mathrm{i}} - f_{\mathrm{pf}})
  \label{eq:Enoe4}
\end{equation}

\textbf{Region 5}---above unity loop gain ($f > f_{\mathrm{i}}$), gain
$\approx f_{\mathrm{c}} / f$:
\begin{equation}
  E_{\mathrm{noe},5}^2
  = e_{\mathrm{nif}}^2 \cdot \frac{f_{\mathrm{c}}^2}{f_{\mathrm{i}}}
  \label{eq:Enoe5}
\end{equation}
Note: \eqref{eq:Enoe5} arises from
$\int_{f_{\mathrm{i}}}^{\infty} (f_{\mathrm{c}}/f)^2 \dd{f}
= f_{\mathrm{c}}^2 / f_{\mathrm{i}}$.

The total output-referred opamp voltage noise is:
\begin{equation}
  E_{\mathrm{noe}}
  = \sqrt{E_{\mathrm{noe},1}^2 + E_{\mathrm{noe},2}^2
          + E_{\mathrm{noe},3}^2 + E_{\mathrm{noe},4}^2
          + E_{\mathrm{noe},5}^2}
  \label{eq:Enoe-total}
\end{equation}

\subsection{Feedback resistor thermal noise}

The Johnson noise of $R_{\mathrm{f}}$ produces a current noise density
$\sqrt{4k_{\mathrm{B}}T/R_{\mathrm{f}}}$, which flows through
$R_{\mathrm{f}}$ to produce output voltage noise.  Since this noise
source sees the same transimpedance transfer function as the signal, the
NEB (\eqref{eq:neb-2nd}) is the appropriate bandwidth:
\begin{equation}
  E_{\mathrm{noR}}
  = \sqrt{4 k_{\mathrm{B}} T R_{\mathrm{f}} \cdot \mathrm{NEB}}
  \label{eq:EnoR}
\end{equation}

\subsection{Shot noise}

Shot noise from DC currents (bias current $I_{\mathrm{b}}$, dark
current $I_{\mathrm{d}}$, and signal photocurrent $I_{\mathrm{p}}$)
produces white current noise with spectral density $\sqrt{2qI}$.  Each
is converted to output voltage noise by the transimpedance
$R_{\mathrm{f}}$ over the NEB:
\begin{equation}
  E_{\mathrm{noi,bias}}
  = R_{\mathrm{f}} \sqrt{2q\,I_{\mathrm{b}} \cdot \mathrm{NEB}}
  \label{eq:Enoi-bias}
\end{equation}
\begin{equation}
  E_{\mathrm{noi,dark}}
  = R_{\mathrm{f}} \sqrt{2q\,I_{\mathrm{d}} \cdot \mathrm{NEB}}
  \label{eq:Enoi-dark}
\end{equation}
\begin{equation}
  E_{\mathrm{noi,sig}}
  = R_{\mathrm{f}} \sqrt{2q\,I_{\mathrm{p}} \cdot \mathrm{NEB}}
  \label{eq:Enoi-sig}
\end{equation}

\subsection{Opamp current noise}

The opamp input current noise floor $i_{\mathrm{nif}}$ (white spectral
density) also flows through $R_{\mathrm{f}}$:
\begin{equation}
  E_{\mathrm{noi,nif}}
  = R_{\mathrm{f}} \, i_{\mathrm{nif}} \sqrt{\mathrm{NEB}}
  \label{eq:Enoi-nif}
\end{equation}

\subsection{Total noise and signal-to-noise ratio}

The total background noise (all sources except signal shot noise) is:
\begin{equation}
  E_{\mathrm{no,bg}}
  = \sqrt{E_{\mathrm{noe}}^2 + E_{\mathrm{noR}}^2
          + E_{\mathrm{noi,bias}}^2 + E_{\mathrm{noi,dark}}^2
          + E_{\mathrm{noi,nif}}^2}
  \label{eq:Eno-bg}
\end{equation}

The signal voltage is
$V_{\mathrm{sig}} = R_{\mathrm{f}} \cdot I_{\mathrm{p}}
= R_{\mathrm{f}} \cdot r_\phi \cdot P_{\mathrm{opt}}$, where
$r_\phi$ is the photodiode responsivity and $P_{\mathrm{opt}}$ is the
incident optical power.  The signal-to-noise ratio is:
\begin{equation}
  \mathrm{SNR}
  = \frac{V_{\mathrm{sig}}}
         {\sqrt{E_{\mathrm{no,bg}}^2 + E_{\mathrm{noi,sig}}^2}}
  \label{eq:snr}
\end{equation}

The ratio $E_{\mathrm{no,bg}} / E_{\mathrm{noi,sig}}$ indicates
whether the system is shot-noise limited (ratio $< 1$) or
background-noise limited (ratio $> 1$).

%======================================================================
\section{Implementation mapping}
\label{sec:implementation}
%======================================================================

\begin{longtable}{@{} p{0.35\textwidth} p{0.40\textwidth}
                      p{0.17\textwidth} @{}}
\toprule
Theory & Implementation & File \\
\midrule
\endfirsthead
\toprule
Theory & Implementation & File \\
\midrule
\endhead
\bottomrule
\endlastfoot
Opamp model, \eqref{eq:opamp-integrator}
  & \texttt{Opamp} dataclass (\texttt{f\_c} field)
  & \texttt{tia\_design.py} \\
Input capacitance $C_{\mathrm{i}}$
  & \texttt{C\_i = C\_d + opamp.C\_icm + opamp.C\_id}
  & \texttt{tia\_design.py} \\
Butterworth $C_{\mathrm{f}}$, \eqref{eq:Cf-exact}
  & \texttt{C\_f = (C\_c * (2*(C\_i+C\_s) - C\_c))**0.5}
  & \texttt{tia\_design.py} \\
Bandwidth, \eqref{eq:f3dB-circuit}
  & \texttt{BW\_t = (f\_c / (2*pi*R\_f*C\_tot))**0.5}
  & \texttt{tia\_design.py} \\
NEB, \eqref{eq:neb-2nd}
  & \texttt{NEB = (pi*sqrt(2)/4) * BW\_t}
  & \texttt{tia\_design.py} \\
Noise gain frequencies, \eqref{eq:fzf}--\eqref{eq:fi}
  & \texttt{f\_zf}, \texttt{f\_pf}, \texttt{f\_i}
  & \texttt{tia\_design.py} \\
Piecewise voltage noise, \eqref{eq:Enoe1}--\eqref{eq:Enoe-total}
  & \texttt{E\_noe1} \ldots\ \texttt{E\_noe5}, \texttt{E\_noe}
  & \texttt{tia\_design.py} \\
Thermal noise, \eqref{eq:EnoR}
  & \texttt{E\_noR}
  & \texttt{tia\_design.py} \\
Shot noise, \eqref{eq:Enoi-bias}--\eqref{eq:Enoi-sig}
  & \texttt{E\_noi\_bias}, \texttt{E\_noi\_dark}, \texttt{E\_noi\_sig}
  & \texttt{tia\_design.py} \\
Current noise, \eqref{eq:Enoi-nif}
  & \texttt{E\_noi\_nif}
  & \texttt{tia\_design.py} \\
SNR, \eqref{eq:snr}
  & \texttt{snr}
  & \texttt{tia\_design.py} \\
Photodiode $C_{\mathrm{d}}(V_{\mathrm{b}})$ model
  & \texttt{Photodiode.C\_d(Vb)}
  & \texttt{tia\_design.py} \\
Multi-opamp comparison
  & \texttt{compare\_opamps(\ldots)}
  & \texttt{tia\_design.py} \\
\end{longtable}

%======================================================================
\section{References}
\label{sec:references}
%======================================================================

\begin{enumerate}
  \item J.~Graeme, \textit{Photodiode Amplifiers: Op Amp Solutions}.
    McGraw-Hill, 1996.
  \item P.~C.~D.~Hobbs,
    \textit{Building Electro-Optical Systems: Making It All Work},
    2nd~ed.\ Wiley, 2009.
\end{enumerate}

\subsection*{Acknowledgment}

This document was prepared with the assistance of Claude (Anthropic).
The author assumes full responsibility for the content.

\end{document}
